\chapter{Introduction}
\label{ch:introduction}

% Core Research Question:
% Why is closed-loop transcranial neurostimulation for epilepsy necessary, why is
% it fundamentally difficult from a control perspective, and what does this thesis
% propose to solve it?

% Contents:
% - Clinical Context: Drug-resistant focal epilepsy and the limitations of resective surgery and pharmacotherapy.
% - Existing Stimulation Paradigms: Open-loop tES and reactive threshold-based triggers (habituation, excessive charge, late intervention).
% - Vision: Receding-horizon optimal control (MPC) to anticipate recruitment and steer neural dynamics under safety budgets.
% - Control Challenges: High-dimensional, stochastic, delayed brain dynamics; severe partial observability (scalp EEG only); inability to evaluate biophysical models online.
% - Proposed Architecture & Key Contributions: End-to-end data-driven MPC pipeline using spectral surrogate Predictors and receding-horizon optimization.
% - Thesis Roadmap: Overview of the narrative and structure of Chapters 2 through 6.
